\documentclass[11pt, a4paper]{article}
\usepackage[utf8]{inputenc}
\usepackage[margin=2.5cm]{geometry}
\usepackage{amsmath, amssymb}
\usepackage{hyperref}

\hypersetup{
    colorlinks=true,
    linkcolor=black,
    citecolor=black,
    urlcolor=blue
}

\title{\textbf{Multi-agentic implementation of a 3D topology optimization framework for compliance and linearized buckling via weighted sum method}}
\author{Gabriele Ricatto}
\date{October 2026}

\begin{document}

\maketitle

\begin{abstract}
This paper presents a rigorous theoretical and computational framework for three-dimensional (3D) multi-objective continuum topology optimization, focusing on the trade-off between structural compliance (global stiffness) and linear buckling stability. A custom Finite Element Method (FEM) solver based on 8-node trilinear hexahedral isoparametric elements (H8) applied to a structured Cartesian voxel grid is implemented. Material distribution is parameterized using the Solid Isotropic Material with Penalization (SIMP) approach, grounded in the 3D Hashin-Shtrikman composite bounds ($p_K = 3$). 

To explore the Pareto-optimal frontier between global stiffness and the fundamental Buckling Load Factor (BLF), a weighted sum method with dynamic $L_1$-norm gradient normalization is developed, guaranteeing Pareto-proper efficiency under Geoffrion's theorem. To eliminate runtime numerical quadrature during the optimization loop, the 3D geometric stiffness matrix is synthesized by decomposing the Cauchy stress tensor into six constant invariant $24 \times 24$ basis matrices, accelerating assembly by two orders of magnitude while suppressing spurious void buckling modes through differential penalization ($p_G \ge p_K$). Exact analytical sensitivities are derived via the adjoint state method requiring a single back-substitution. The numerical architecture is optimized for execution on commodity hardware with 8GB RAM (AMD Ryzen 7 4700U), consuming under $100\,\text{MB}$ peak RAM for standard test domains, and seamlessly interfaces with a watertight STL surface extraction pipeline.
\end{abstract}

\section{Introduction}
In conventional minimum compliance topology optimization governed purely by linear elasticity, applying an external force $\mathbf{F}$ or its exact inverse $-\mathbf{F}$ produces an identical optimized topology. Static equilibrium is governed by Hooke's law:
\begin{equation}
\mathbf{K}(\boldsymbol{\rho}) \mathbf{U} = \mathbf{F}
\end{equation}
where $\mathbf{K}$ is the symmetric, positive-definite global stiffness matrix, $\mathbf{U}$ is the nodal displacement vector, and $\mathbf{F}$ is the applied external force vector. Inverting the load vector to $-\mathbf{F}$ yields an identically reversed displacement field:
\begin{equation}
\mathbf{U}_{-\mathbf{F}} = \mathbf{K}^{-1}(-\mathbf{F}) = -\mathbf{U}
\end{equation}
Structural compliance $C$ represents the total external work done by the loads:
\begin{equation}
C = \mathbf{F}^T \mathbf{U} = \mathbf{U}^T \mathbf{K} \mathbf{U} = \sum_{e=1}^m E(\rho_e) \mathbf{u}_e^T \mathbf{k}_0 \mathbf{u}_e
\end{equation}
Evaluating the compliance under the reversed load $-\mathbf{F}$ gives:
\begin{equation}
C_{-\mathbf{F}} = (-\mathbf{F})^T (-\mathbf{U}) = \mathbf{F}^T \mathbf{U} = C
\end{equation}
Because the strain energy is a quadratic form in the nodal displacements, it remains strictly positive and insensitive to the sign of the displacement vector. Under small-strain linear elasticity and symmetric material behavior ($E_{\text{tension}} = E_{\text{compression}}$), tensile and compressive stresses contribute identically to structural compliance. 

Consequently, a compliance-only optimizer cannot differentiate between a slender tension tie and an Euler column subjected to axial compression. Under compression, slender members undergo catastrophic geometric instability (buckling) at loads orders of magnitude below their yield limit. Breaking this artificial symmetry requires the explicit inclusion of Linearized Buckling Analysis (LBA) into the optimization objective [1, 7].

\section{FEM Implementation: 3D Isoparametric H8 Element}
For the implementation of the 3D FEM solver, we rely on eight-node trilinear hexahedral isoparametric elements (H8) applied to a fully structured Cartesian voxel grid. The rectangular elements have dimensions $dx \times dy \times dz$ and volume $V_e = dx \cdot dy \cdot dz$, aligned with the global reference system, with the origin of the natural coordinates $(\xi, \eta, \zeta) \in [-1, 1]^3$ centered in the centroid of each element $(x_c, y_c, z_c)$:
\begin{equation}
x = x_c + \frac{dx}{2}\xi, \quad y = y_c + \frac{dy}{2}\eta, \quad z = z_c + \frac{dz}{2}\zeta
\end{equation}
with $\xi, \eta, \zeta \in [-1, 1]$. The Jacobian matrix $[\mathbf{J}]$ is strictly diagonal:
\begin{equation}
[\mathbf{J}] = \begin{bmatrix}
\frac{\partial x}{\partial \xi} & \frac{\partial y}{\partial \xi} & \frac{\partial z}{\partial \xi} \\[4pt]
\frac{\partial x}{\partial \eta} & \frac{\partial y}{\partial \eta} & \frac{\partial z}{\partial \eta} \\[4pt]
\frac{\partial x}{\partial \zeta} & \frac{\partial y}{\partial \zeta} & \frac{\partial z}{\partial \zeta}
\end{bmatrix} = \begin{bmatrix}
dx/2 & 0 & 0 \\
0 & dy/2 & 0 \\
0 & 0 & dz/2
\end{bmatrix} \implies \det(\mathbf{J}) = \frac{1}{8} dx \, dy \, dz = \frac{1}{8} V_e
\end{equation}
and its inverse is:
\begin{equation}
[\mathbf{J}]^{-1} = \begin{bmatrix}
2/dx & 0 & 0 \\
0 & 2/dy & 0 \\
0 & 0 & 2/dz
\end{bmatrix}
\end{equation}

The trilinear shape functions $N_i(\xi, \eta, \zeta)$ for the 8 local nodes are:
\begin{equation}
N_i(\xi, \eta, \zeta) = \frac{1}{8}(1 + \xi_i \xi)(1 + \eta_i \eta)(1 + \zeta_i \zeta), \quad i \in \{1, 2, \dots, 8\}
\end{equation}
where $(\xi_i, \eta_i, \zeta_i) \in \{-1, 1\}^3$ denote the corner coordinates in natural space.

\subsection{Kinematic Degrees of Freedom and B-Matrix}
Each node possesses 3 translational degrees of freedom $(u, v, w)$, yielding $8 \times 3 = 24$ DOFs per element. The element displacement vector is:
\begin{equation}
\mathbf{u}_e = \begin{bmatrix} u_1 & v_1 & w_1 & u_2 & v_2 & w_2 & \dots & u_8 & v_8 & w_8 \end{bmatrix}^T \in \mathbb{R}^{24}
\end{equation}
The engineering strain vector $\boldsymbol{\varepsilon} = [\varepsilon_{xx}, \varepsilon_{yy}, \varepsilon_{zz}, \gamma_{yz}, \gamma_{zx}, \gamma_{xy}]^T \in \mathbb{R}^6$ is interpolated via the $6 \times 24$ strain-displacement matrix $\mathbf{B}$:
\begin{equation}
\boldsymbol{\varepsilon}(\xi, \eta, \zeta) = \mathbf{B}(\xi, \eta, \zeta) \mathbf{u}_e = \begin{bmatrix} \mathbf{B}_1 & \mathbf{B}_2 & \dots & \mathbf{B}_8 \end{bmatrix} \mathbf{u}_e
\end{equation}
where each $6 \times 3$ nodal submatrix $\mathbf{B}_i$ is defined as:
\begin{equation}
\mathbf{B}_i = \begin{bmatrix}
\frac{2}{dx} \frac{\partial N_i}{\partial \xi} & 0 & 0 \\[3pt]
0 & \frac{2}{dy} \frac{\partial N_i}{\partial \eta} & 0 \\[3pt]
0 & 0 & \frac{2}{dz} \frac{\partial N_i}{\partial \zeta} \\[3pt]
0 & \frac{2}{dz} \frac{\partial N_i}{\partial \zeta} & \frac{2}{dy} \frac{\partial N_i}{\partial \eta} \\[3pt]
\frac{2}{dz} \frac{\partial N_i}{\partial \zeta} & 0 & \frac{2}{dx} \frac{\partial N_i}{\partial \xi} \\[3pt]
\frac{2}{dy} \frac{\partial N_i}{\partial \eta} & \frac{2}{dx} \frac{\partial N_i}{\partial \xi} & 0
\end{bmatrix}
\end{equation}

\subsection{Element Stiffness Matrix \texorpdfstring{$\mathbf{k}_0$}{k\_0} ($24 \times 24$)}
For isotropic 3D linear elasticity characterized by Young's modulus $E_0$ and Poisson's ratio $\nu$, the constitutive matrix $\mathbf{D} \in \mathbb{R}^{6 \times 6}$ is:
\begin{equation}
\mathbf{D} = \frac{E_0}{(1+\nu)(1-2\nu)} \begin{bmatrix}
1-\nu & \nu & \nu & 0 & 0 & 0 \\
\nu & 1-\nu & \nu & 0 & 0 & 0 \\
\nu & \nu & 1-\nu & 0 & 0 & 0 \\
0 & 0 & 0 & \frac{1-2\nu}{2} & 0 & 0 \\
0 & 0 & 0 & 0 & \frac{1-2\nu}{2} & 0 \\
0 & 0 & 0 & 0 & 0 & \frac{1-2\nu}{2}
\end{bmatrix}
\end{equation}
The base element stiffness matrix $\mathbf{k}_0 \in \mathbb{R}^{24 \times 24}$ is evaluated via exact $2 \times 2 \times 2$ Gauss-Legendre quadrature:
\begin{equation}
\mathbf{k}_0 = \int_{-1}^1 \int_{-1}^1 \int_{-1}^1 \mathbf{B}^T \mathbf{D} \mathbf{B} \det(\mathbf{J}) \, d\xi \, d\eta \, d\zeta
\end{equation}
Spectral analysis confirms that $\mathbf{k}_0$ possesses \textbf{exactly 6 zero eigenvalues} (three rigid translations and three rigid rotations) and \textbf{18 strictly positive eigenvalues} corresponding to independent elastic strain energy modes, verifying freedom from hourglassing instabilities.

\subsection{3D Geometric Stiffness Decomposition into 6 Invariant Basis Matrices}
The elemental geometric stiffness matrix $\mathbf{G}_0^e$ is derived from the linearization of the Principle of Virtual Work under finite Green-Lagrange strain:
\begin{equation}
\mathbf{G}_0^e = \int_{\Omega_e} \mathbf{B}_{\text{nl}}^T \mathbf{M}_\sigma \mathbf{B}_{\text{nl}} \, d\Omega
\end{equation}
where $\mathbf{B}_{\text{nl}} \in \mathbb{R}^{9 \times 24}$ contains the spatial gradients of the displacement components, and $\mathbf{M}_\sigma = \mathbf{I}_3 \otimes \boldsymbol{\sigma}_e \in \mathbb{R}^{9 \times 9}$ embeds the $3 \times 3$ Cauchy stress tensor:
\begin{equation}
\boldsymbol{\sigma}_e = \begin{bmatrix}
\sigma_{xx} & \tau_{xy} & \tau_{xz} \\
\tau_{xy} & \sigma_{yy} & \tau_{yz} \\
\tau_{xz} & \tau_{yz} & \sigma_{zz}
\end{bmatrix}
\end{equation}
Decomposing $\mathbf{M}_\sigma$ into its 6 independent scalar stress components [7]:
\begin{equation}
\mathbf{M}_\sigma = \sigma_{xx} \mathbf{T}_{xx} + \sigma_{yy} \mathbf{T}_{yy} + \sigma_{zz} \mathbf{T}_{zz} + \tau_{yz} \mathbf{T}_{yz} + \tau_{xz} \mathbf{T}_{xz} + \tau_{xy} \mathbf{T}_{xy}
\end{equation}
where $\mathbf{T}_k = \mathbf{I}_3 \otimes \mathbf{t}_k$ are constant binary basis matrices. Substituting into the integral:
\begin{equation}
\mathbf{G}_0^e = \sum_{k \in \{xx, yy, zz, yz, xz, xy\}} \sigma_{k, e} \mathbf{G}_{0, k}
\end{equation}
where the six constant $24 \times 24$ geometric stiffness basis matrices are defined as:
\begin{equation}
\mathbf{G}_{0, k} = \int_{\Omega_e} \mathbf{B}_{\text{nl}}^T \mathbf{T}_k \mathbf{B}_{\text{nl}} \, d\Omega = V_e \cdot \mathbf{B}_{\text{nl}}(0,0,0)^T \mathbf{T}_k \mathbf{B}_{\text{nl}}(0,0,0)
\end{equation}
\textbf{Computational Advantage:} Because all six basis matrices $\mathbf{G}_{0, k}$ depend solely on element dimensions $(dx, dy, dz)$, they are precomputed \textbf{exactly once} during initialization. Inside the optimization iterations, element geometric stiffness is formed by simple scalar multiplication, accelerating assembly by more than two orders of magnitude.

\section{Principal Stress Decomposition: Signed Stress Metric for Tension and Compression}
At each hexahedral element $e$, the centroidal engineering strain $\boldsymbol{\varepsilon}_e = \mathbf{B}(0,0,0) \mathbf{u}_e$ yields the elemental Cauchy stress tensor $\boldsymbol{\sigma}_e \in \mathbb{R}^{3 \times 3}$:
\begin{equation}
\boldsymbol{\sigma}_e = \begin{bmatrix}
\sigma_{xx} & \tau_{xy} & \tau_{xz} \\
\tau_{xy} & \sigma_{yy} & \tau_{yz} \\
\tau_{xz} & \tau_{yz} & \sigma_{zz}
\end{bmatrix}
\end{equation}
Because $\boldsymbol{\sigma}_e$ is real and symmetric ($\boldsymbol{\sigma}_e = \boldsymbol{\sigma}_e^T$), the spectral theorem guarantees the existence of three real eigenvalues and a mutually orthogonal triad of principal unit vectors $\{\mathbf{n}_I, \mathbf{n}_{II}, \mathbf{n}_{III}\}$. The stress tensor admits the spectral decomposition:
\begin{equation}
\boldsymbol{\sigma}_e = \sum_{p \in \{I, II, III\}} \sigma_p \, \mathbf{n}_p \otimes \mathbf{n}_p = \mathbf{Q} \boldsymbol{\Lambda} \mathbf{Q}^T
\end{equation}
where $\mathbf{Q} = [\mathbf{n}_I, \mathbf{n}_{II}, \mathbf{n}_{III}] \in \mathrm{SO}(3)$ is the orthogonal transformation matrix, and the diagonal matrix is $\boldsymbol{\Lambda} = \operatorname{diag}(\sigma_I, \sigma_{II}, \sigma_{III})$. The principal stresses are algebraically ordered:
\begin{equation}
\sigma_I \ge \sigma_{II} \ge \sigma_{III}
\end{equation}
where $\sigma_I$ denotes the maximum principal tensile stress, $\sigma_{III}$ denotes the maximum compressive stress (the most negative eigenvalue), and $\sigma_{II}$ represents the intermediate principal stress.

In classical structural analysis, scalar stress invariants such as the von Mises equivalent stress $\sigma_{\text{vM}}$:
\begin{equation}
\sigma_{\text{vM}} = \sqrt{\frac{1}{2} \left[ (\sigma_I - \sigma_{II})^2 + (\sigma_{II} - \sigma_{III})^2 + (\sigma_{III} - \sigma_I)^2 \right]} \ge 0
\end{equation}
are strictly non-negative and blind to hydrostatic sign reversals. As a result, $\sigma_{\text{vM}}$ cannot distinguish between an intense tensile state and a severe compressive state.

To resolve this limitation and distinguish tensile load paths from stability-critical members, we define the \textit{Signed Principal Stress} metric $\sigma_{\text{signed}}$:
\begin{equation}
\sigma_{\text{signed}} = \begin{cases} 
\sigma_I & \text{if } |\sigma_I| \ge |\sigma_{III}| \\[6pt]
\sigma_{III} & \text{if } |\sigma_{III}| > |\sigma_I| 
\end{cases}
\end{equation}
This scalar field isolates the dominant principal stress magnitude while preserving its physical sign:
\begin{itemize}
    \item \textbf{Tensile Chords ($\sigma_{\text{signed}} > 0$):} Elements where $\sigma_{\text{signed}} > 0$ reside in a tensile stress regime governed by $\sigma_I$. Tensile forces induce positive-definite geometric stiffness contributions (stress stiffening), rendering tensile ties inherently stable against geometric bifurcation. Under combined compliance and stability optimization, these members can be realized as highly slender, direct tensile ties without incurring buckling penalties.
    \item \textbf{Euler Buckling-Prone Compressive Chords ($\sigma_{\text{signed}} < 0$):} Elements where $\sigma_{\text{signed}} < 0$ are dominated by compressive stress ($\sigma_{III} < 0$). In accordance with Euler's critical buckling load formulation:
    \begin{equation}
    P_{\text{cr}} = \frac{\pi^2 E I}{(K L)^2}
    \end{equation}
    slender structural members subjected to compressive axial loads are prone to catastrophic elastic bifurcation instabilities at stress levels well below the material yield limit. The compressive eigenvalues in $\boldsymbol{\sigma}_e$ generate negative contributions to the geometric stiffness matrix $\mathbf{G}$, driving the structure toward geometric collapse ($\det(\mathbf{K} + \lambda \mathbf{G}) \to 0$).
\end{itemize}
By monitoring $\sigma_{\text{signed}}$ across the design domain, the topology optimization framework explicitly identifies compressive struts that require larger second moments of area $I$—motivating the formation of lateral bracing, webbed cross-sections, or thickened compressive arches—while permitting pure tension ties to remain slender.

\section{Optimization Problem Setup}
The spatial domain is discretized into $m$ equal elements. The pseudo-density field is defined as $\hat{\mathbf{x}} = \{\hat{x}_e\}_{e=1}^m \in [0, 1]^m$. Intermediate densities are filtered through a 3D spherical cone filter ($r_{\min}$) [2, 7]:
\begin{equation}
\tilde{x}_e = \frac{\sum_{i=1}^m H_{ei} x_i}{\sum_{i=1}^m H_{ei}}, \quad H_{ei} = \max(0, \, r_{\min} - \text{dist}(\Omega_e, \Omega_i))
\end{equation}
and projected toward discrete $0/1$ boundaries via smoothed Heaviside projection [8]:
\begin{equation}
\hat{x}_e = \frac{\tanh(\beta \eta) + \tanh(\beta (\tilde{x}_e - \eta))}{\tanh(\beta \eta) + \tanh(\beta (1 - \eta))}
\end{equation}
with threshold $\eta = 0.5$ and continuation parameter $\beta \in [1, 32]$.

\subsection{SIMP Interpolation and Differential Penalization}
The elastic and geometric stiffnesses are parameterized using differential penalization [1, 7]:
\begin{align}
E_K(\hat{x}_e) &= E_{\min} + (E_0 - E_{\min}) \hat{x}_e^{p_K} \\
E_G(\hat{x}_e) &= E_0 \hat{x}_e^{p_G}
\end{align}
with $p_K = 3.0$ (grounded in 3D Hashin-Shtrikman bounds) and $p_G = 3.0$ (or $6.0$). In void regions ($\hat{x}_e \to 0$):
\begin{equation}
\frac{E_G(\hat{x}_e)}{E_K(\hat{x}_e)} \xrightarrow{\hat{x}_e \to 0} \frac{0}{E_{\min}} = 0
\end{equation}
This completely suppresses spurious void buckling modes in empty space without corrupting the physical low-frequency structural spectrum.

\subsection{Linearized Buckling Eigenvalue Problem}
Linearized structural stability is governed by the generalized eigenvalue problem:
\begin{equation}
(\mathbf{K} + \lambda_j \mathbf{G}) \boldsymbol{\phi}_j = \mathbf{0}
\end{equation}
To prevent singularities in void elements, the reciprocal eigenvalue $\mu_j = 1/\lambda_j$ is evaluated [1, 7]:
\begin{equation}
(\mathbf{G} + \mu_j \mathbf{K}) \boldsymbol{\phi}_j = \mathbf{0}
\end{equation}
where the fundamental critical buckling load factor is $\lambda_1 = 1/\mu_1$ with $\mu_1 = \max_j(\mu_j)$.

\subsection{Multi-Objective Formulation and Dynamic \texorpdfstring{$L_1$}{L1} Sensitivity Scaling}
To ensure global static load bearing while simultaneously stiffening compressive members against buckling, we formulate the multi-objective optimization problem:
\begin{equation}
\min_{\mathbf{x}} \left\{ C(\hat{\mathbf{x}}), \, -\mu_1(\hat{\mathbf{x}}) \right\} \quad \text{subject to} \quad \frac{1}{m}\sum_{e=1}^m \hat{x}_e \le V_f
\label{eq:multiobj_prob}
\end{equation}
Because these two objectives are competing, the solution space forms a Pareto-optimal frontier. We scalarize the problem into a combined objective function $F(\hat{\mathbf{x}})$ using the weighted sum method [3]:
\begin{equation}
F(\hat{\mathbf{x}}) = (1 - \alpha) C + \alpha \mu_1, \quad \alpha \in [0, 1]
\label{eq:scalarized_obj}
\end{equation}
Because compliance and eigenvalue sensitivities differ by orders of magnitude, we implement adaptive dynamic sensitivity scaling. At each iteration $i$, scaling factors $S_C^{(i)}$ and $S_\mu^{(i)}$ are evaluated as the mean absolute values ($L_1$-norms) of the raw gradients:
\begin{equation}
S_C^{(i)} = \frac{1}{m}\sum_{e=1}^m \left| \frac{\partial C^{(i)}}{\partial \hat{x}_e} \right|, \quad 
S_\mu^{(i)} = \frac{1}{m}\sum_{e=1}^m \left| \frac{\partial \mu_1^{(i)}}{\partial \hat{x}_e} \right|
\end{equation}
The aggregated sensitivity passed to the optimizer is:
\begin{equation}
\frac{\partial F^{(i)}}{\partial \hat{x}_e} = (1 - \alpha) \frac{1}{S_C^{(i)}} \frac{\partial C^{(i)}}{\partial \hat{x}_e} + \alpha \frac{1}{S_\mu^{(i)}} \frac{\partial \mu_1^{(i)}}{\partial \hat{x}_e}
\end{equation}
Under Geoffrion's theorem [10], dividing by positive scaling norms mathematically guarantees convergence to a properly efficient Pareto-optimal solution.

\section{Accelerated Heaviside Continuation for Slender 3D Topologies}
To transform intermediate densities into distinct $\{0, 1\}$ structural topologies, filtered densities $\tilde{x}_e$ are mapped through a smoothed Heaviside projection [8]:
\begin{equation}
\hat{x}_e = \frac{\tanh(\beta \eta) + \tanh(\beta (\tilde{x}_e - \eta))}{\tanh(\beta \eta) + \tanh(\beta (1 - \eta))}
\end{equation}
where $\eta \in (0, 1)$ is the threshold parameter (typically $\eta = 0.5$) and $\beta \in [1, \infty)$ dictates the sharpness of the transition interface.

\subsection{Pathologies of Static Low \texorpdfstring{$\beta$}{beta} and Abrupt High \texorpdfstring{$\beta$}{beta}}
The magnitude of $\beta$ governs the regularization of the step function. When $\beta$ is maintained at a low constant value ($\beta = 1$), the projection remains smooth and well-behaved, facilitating broad topological exploration. However, persistent low $\beta$ produces wide ``grey zones'' of intermediate density ($0 < \hat{x}_e < 1$) across element boundaries. In 3D linearized buckling optimization, extended grey zones distort the differential penalization ratio $E_G(\hat{x}_e) / E_K(\hat{x}_e)$, creating non-physical boundary compliance and obscuring localized structural bifurcation modes.

Conversely, abruptly initializing or stepping to high $\beta$ (e.g., $\beta \ge 16$) causes severe gradient vanishing away from the threshold interface $\eta$:
\begin{equation}
\frac{\partial \hat{x}_e}{\partial \tilde{x}_e} = \frac{\beta \operatorname{sech}^2(\beta(\tilde{x}_e - \eta))}{\tanh(\beta \eta) + \tanh(\beta(1 - \eta))} \xrightarrow{|\tilde{x}_e - \eta| > \delta} 0
\end{equation}
This premature vanishing traps the design variables in sub-optimal local extrema, freezing chunky, stubby, and over-bulked structural members that fail to branch into slender, high-efficiency truss topologies.

\subsection{Accelerated Continuation Formulation}
To eliminate intermediate grey zones while simultaneously resolving slender structural features, an accelerated continuation schedule for $\beta$ is enforced:
\begin{equation}
\beta^{(k+1)} = \min\left( \beta_{\max}, \; \gamma_\beta \cdot \beta^{(k)} \right)
\end{equation}
where $\beta^{(0)} = 1$, $\gamma_\beta \in [1.5, 2.0]$ is the continuation multiplier (e.g., geometric doubling $\beta \in \{1, 2, 4, 8, 16, 32, 64\}$), and $\beta_{\max} = 32$ (or $64$). The continuation step is updated according to:
\begin{equation}
\text{Update } \beta \quad \text{if} \quad \left( k \pmod{\Delta N_\beta} = 0 \right) \quad \text{or} \quad \left( \|\mathbf{x}^{(k)} - \mathbf{x}^{(k - \Delta k)}\|_\infty < \epsilon_\beta \right)
\end{equation}
with update cadence $\Delta N_\beta \in [8, 12]$ iterations and convergence tolerance $\epsilon_\beta = 0.01$.

During the initial phase ($\beta \le 4$), low projection steepness permits the continuous material to flow globally, establishing optimal structural connectivity between supports and loads. As the continuation accelerates $\beta$ into higher regimes ($\beta \ge 8$), the transition layer narrows sharply to the element width, forcing intermediate voxels strictly to 0 or 1. This rapid boundary sharpening eliminates grey zones and prevents the accumulation of chunky, stubby material blocks, compelling the optimizer to discover slender, high-aspect-ratio tensile ties and optimized triangulated compressive bracing.

\section{Sensitivity Analysis}
The compliance derivative is evaluated as [7]:
\begin{equation}
\frac{\partial C}{\partial \hat{x}_e} = - \mathbf{u}_e^T \mathbf{k}_0 \mathbf{u}_e \frac{\partial E_K(\hat{x}_e)}{\partial \hat{x}_e} = - p_K \hat{x}_e^{p_K - 1}(E_0 - E_{\min}) \mathbf{u}_e^T \mathbf{k}_0 \mathbf{u}_e
\end{equation}

For the fundamental buckling reciprocal eigenvalue $\mu_1$, the exact sensitivity is derived via the adjoint state method [4, 7]:
\begin{equation}
\frac{\partial \mu_1}{\partial \hat{x}_e} = - \left[ \boldsymbol{\phi}_{1, e}^T \frac{\partial \mathbf{G}_e}{\partial \hat{x}_e} \boldsymbol{\phi}_{1, e} + \mu_1 \boldsymbol{\phi}_{1, e}^T \frac{\partial \mathbf{k}_e}{\partial \hat{x}_e} \boldsymbol{\phi}_{1, e} - \mathbf{w}_{1, e}^T \frac{\partial \mathbf{k}_e}{\partial \hat{x}_e} \mathbf{u}_e \right]
\end{equation}
where:
\begin{itemize}
    \item $\boldsymbol{\phi}_1$: eigenvector / fundamental buckling mode shape, normalized such that $\boldsymbol{\phi}_1^T \mathbf{K} \boldsymbol{\phi}_1 = 1$.
    \item $\mu_1$: fundamental reciprocal eigenvalue ($1/\lambda_1$).
    \item $\mathbf{u}_e$: elemental static displacement vector under external loads $\mathbf{F}$.
    \item $\mathbf{w}_1$: global adjoint displacement vector, obtained by solving:
    \begin{equation}
    \mathbf{K} \mathbf{w}_1 = \mathbf{F}_{\text{adj}} = \sum_{e=1}^m E_G(\hat{x}_e) \mathbf{B}_0^T \mathbf{D}^T \mathbf{P}_e
    \end{equation}
    where $P_{e, k} = \boldsymbol{\phi}_{1, e}^T \mathbf{G}_{0, k} \boldsymbol{\phi}_{1, e}$ for $k \in \{xx, yy, zz, yz, xz, xy\}$.
\end{itemize}
Crucially, evaluating $\mathbf{w}_1$ requires \textbf{zero new matrix factorizations}: the pre-factored stiffness matrix from the forward static solve $\mathbf{K} \mathbf{U} = \mathbf{F}$ is directly reused for a single rapid back-substitution.

\section{Updating Process}
The physical sensitivities are back-propagated through the Heaviside projection and filter chain rule:
\begin{equation}
\frac{\partial \tilde{F}}{\partial x_i} = \sum_{e=1}^m \frac{H_{ei}}{\sum_k H_{ek}} \left( \frac{\partial F}{\partial \hat{x}_e} \frac{\partial \hat{x}_e}{\partial \tilde{x}_e} \right)
\end{equation}
The design variables are updated using the Optimality Criteria (OC) scheme \cite{sigmund2001, bendsoe2003} (for multi-constrained formulations and thermo-elastic problems with non-monotonic sensitivities, this heuristic update is superseded by the Method of Moving Asymptotes \cite{svanberg1987}, detailed in Section \ref{sec:mma}):
\begin{equation}
x_e^{(k+1)} = \begin{cases}
\max(x_{\min}, \, x_e^{(k)} - \zeta) & \text{if } x_e^{(k)} B_e \le \max(x_{\min}, \, x_e^{(k)} - \zeta) \\[4pt]
x_e^{(k)} B_e & \text{if } \max(x_{\min}, \, x_e^{(k)} - \zeta) < x_e^{(k)} B_e < \min(1, \, x_e^{(k)} + \zeta) \\[4pt]
\min(1, \, x_e^{(k)} + \zeta) & \text{if } \min(1, \, x_e^{(k)} + \zeta) \le x_e^{(k)} B_e
\end{cases}
\end{equation}
where $\zeta = 0.2$ is the iteration move limit, $x_{\min} = 0.001$, and:
\begin{equation}
B_e = \sqrt{\max\left(0, \, -\frac{\partial \tilde{F} / \partial x_e}{\lambda_L \cdot \partial V / \partial x_e}\right)}
\end{equation}
The volume Lagrange multiplier $\lambda_L$ is determined via bisection search until volume conservation $|\frac{1}{m}\sum_e \hat{x}_e - V_f| < 10^{-4}$ is satisfied.

\section{Post-Processing: Discrete Surface Smoothing via Taubin Volume-Preserving Filtering}
Upon convergence of the filtered and projected pseudo-density field $\hat{\mathbf{x}}^*$, an explicit geometric boundary of the structural topology is extracted at the discrete isosurface $\hat{x}_{\text{iso}} = 0.5$. In the Cartesian voxel framework, this is accomplished by traversing the computational grid to detect all external boundary quad facets separating solid voxels ($\hat{x}_e \ge 0.5$) from void voxels ($\hat{x}_e < 0.5$). Each rectangular quadrilateral facet is split along its shortest diagonal into two coplanar triangles, yielding an orientable, watertight 2-manifold triangular boundary mesh:
\begin{equation}
\mathcal{M} = (\mathcal{V}, \mathcal{F})
\end{equation}
parameterized by the vertex coordinate set $\mathcal{V} = \{\mathbf{v}_i \in \mathbb{R}^3\}_{i=1}^{N_v}$ and the triangular face connectivity set $\mathcal{F} = \{f_k\}_{k=1}^{N_f}$.

Because the boundary mesh is extracted directly from orthogonal hexahedral cells, raw exported isosurfaces exhibit severe high-frequency voxel staircase artifacts. These microscopic discontinuities create spurious stress concentrations in downstream validation analyses and introduce unacceptable surface roughness for additive manufacturing and slicer toolpath generation.

\subsection{Discrete Laplace-Beltrami Operator on 2-Manifold Boundary Meshes}
Surface smoothing is formulated directly on the discrete graph structure of $\mathcal{M}$ through the discrete Laplace-Beltrami operator. For any boundary vertex $\mathbf{v}_i \in \mathcal{V}$, the discrete surface Laplacian $\Delta \mathbf{v}_i$ represents the vector displacement between $\mathbf{v}_i$ and the convex combination of its adjacent 1-ring neighbors:
\begin{equation}
\Delta \mathbf{v}_i = \sum_{j \in \mathcal{N}(i)} w_{ij} (\mathbf{v}_j - \mathbf{v}_i)
\end{equation}
where $\mathcal{N}(i) = \{j \in \mathcal{V} : (i, j) \in \mathcal{E}\}$ denotes the set of immediate topological neighbors connected to vertex $i$ by an edge in $\mathcal{E}$, and $w_{ij}$ are non-negative normalized edge weights satisfying the partition of unity:
\begin{equation}
\sum_{j \in \mathcal{N}(i)} w_{ij} = 1
\end{equation}
Under the discrete umbrella operator (uniform weighting), the weights are simply evaluated as:
\begin{equation}
w_{ij} = \frac{1}{|\mathcal{N}(i)|}
\end{equation}
where $|\mathcal{N}(i)|$ is the valence (degree) of vertex $\mathbf{v}_i$.

\subsection{Volumetric Shrinkage of Classical Laplacian Smoothing}
Classical Laplacian smoothing iteratively updates vertex coordinates through the forward Euler discretization of the isotropic surface heat diffusion equation $\partial \mathbf{v} / \partial t = \Delta \mathbf{v}$:
\begin{equation}
\mathbf{v}_i^{(k+1)} = \mathbf{v}_i^{(k)} + \lambda \Delta \mathbf{v}_i^{(k)}, \quad \lambda > 0
\end{equation}
While explicit Laplacian smoothing effectively damps high-frequency spatial noise, it acts as a low-pass Gaussian filter that induces severe global volume loss and geometric shrinkage. For slender structural members—such as thin tensile ties and delicate compressive trusses—repeated Laplacian iterations steadily erode cross-sectional dimensions, substantially shrinking the structural volume and degrading the buckling load capacity.

\subsection{Dual-Step Taubin Volume-Preserving Smoothing Filter}
To suppress high-frequency voxel staircase noise without shrinking the macroscopic volume or rounding slender structural chords, we incorporate the non-shrinking dual-step filter formulated by Taubin \cite{taubin1995}. Within each smoothing cycle $k \in \{0, 1, \dots, N_{\text{cycles}}-1\}$, the algorithm alternates between an inward diffusion step and an outward anti-diffusion step:
\begin{align}
\mathbf{v}_i^{(2k+1)} &= \mathbf{v}_i^{(2k)} + \lambda \Delta \mathbf{v}_i^{(2k)} \\
\mathbf{v}_i^{(2k+2)} &= \mathbf{v}_i^{(2k+1)} + \mu \Delta \mathbf{v}_i^{(2k+1)}
\end{align}
with positive shrink parameter $\lambda > 0$ and negative unshrink parameter $\mu < -\lambda < 0$, satisfying the volume-preservation inequality:
\begin{equation}
\frac{1}{\lambda} + \frac{1}{\mu} < 0
\end{equation}
In the spectral domain of the discrete Laplace-Beltrami operator, with spatial frequency eigenvalues $k \in [0, k_{\max}]$ (where $k_{\max} \le 2$ on normalized meshes), the transfer function over a complete two-step iteration cycle is:
\begin{equation}
f(k) = (1 - \lambda k)(1 - \mu k) = 1 - (\lambda + \mu) k + \lambda \mu k^2
\end{equation}
The mathematical condition $\mu < -\lambda < 0$ and $\frac{1}{\lambda} + \frac{1}{\mu} < 0$ guarantees that:
\begin{enumerate}
    \item \textbf{Zero Low-Frequency Shrinkage:} At low spatial frequencies ($k \to 0$), the transfer function maintains unit gain ($f(0) = 1$) with positive slope $f'(0) = -(\lambda + \mu) > 0$. The anti-diffusive step ($\mu < 0$) counteracts the shrinkage induced by the positive diffusion step ($\lambda > 0$), perfectly stabilizing the global enclosed volume.
    \item \textbf{Attenuation of High-Frequency Staircase Artifacts:} For high spatial frequencies $k \in (k_{\text{pass}}, 2]$, the magnitude of the transfer function satisfies $|f(k)| < 1$, ensuring monotonic damping of the Nyquist-frequency staircase artifacts associated with voxel edge discretization.
\end{enumerate}
In the production pipeline, stable filter parameters $\lambda = 0.50$ and $\mu = -0.53$ are applied over $N_{\text{cycles}} = 10 \text{--} 20$ iterations. This produces smooth, curvature-continuous, and volume-preserving 2-manifold triangular meshes that strictly preserve the mechanical stiffness, cross-sectional area, and Euler buckling resistance of slender 3D structural topologies.

\section{Numerical Performance on 8GB Hardware}
To guarantee execution on standard laptop hardware (AMD Ryzen 7 4700U with 8GB RAM), the solver employs:
\begin{enumerate}
    \item \textbf{Preconditioned Conjugate Gradient (PCG):} Jacobi diagonal preconditioning coupled with displacement warm-starting achieves convergence in $O(\sqrt{\kappa})$ iterations while using only $\approx 12\,\text{KB}$ RAM per element (reducing solver memory by $18\times$ compared to direct LU factorizations).
    \item \textbf{Invariant Assembly Indexing:} Sparse row and column coordinate vectors $(iK, jK)$ are pre-computed once and shared between elastic stiffness $\mathbf{K}$ and geometric stiffness $\mathbf{G}$.
    \item \textbf{Watertight STL Export:} A boundary-face extraction algorithm identifies only external quad facets between solid ($\hat{x}_e \ge \text{threshold}$) and void voxels, producing closed, manifold triangular meshes without internal partition planes, directly consumable by additive manufacturing slicers and PicoGK OpenVDB pipelines.
\end{enumerate}

\begin{table}[h!]
\centering
\footnotesize
\setlength{\tabcolsep}{3pt}
\caption{Memory scaling and performance across 3D voxel grid resolutions (Ryzen 7 4700U, 8GB RAM).}
\vspace{4pt}
\begin{tabular}{|cccccc|}
\hline
\textbf{Grid ($N_x \times N_y \times N_z$)} & \textbf{Elements} & \textbf{Total DOFs} & \textbf{Sparse RAM} & \textbf{Peak RAM} & \textbf{Suitability} \\
\hline
$10 \times 10 \times 10$ & 1,000 & 3,993 & $\sim 7$ MB & \textbf{$<$ 60 MB} & Fast Unit Test \\
$20 \times 10 \times 10$ & 2,000 & 7,623 & $\sim 14$ MB & \textbf{$<$ 120 MB} & Interactive Dashboard \\
$30 \times 15 \times 10$ & 4,500 & 16,368 & $\sim 31$ MB & \textbf{$<$ 310 MB} & High-Fidelity Domain \\
$40 \times 20 \times 10$ & 8,000 & 28,413 & $\sim 55$ MB & \textbf{$<$ 620 MB} & Production Domain \\
\hline
\end{tabular}
\end{table}

\section{Development Log (Sprint 1)}
The following critical issues were resolved during Sprint 1:
\begin{itemize}
    \item \textbf{Fix 1: Zero-Energy Hourglassing Modes.} The geometric stiffness basis matrices $\mathbf{G}_{0,k}$ were originally calculated using centroidal 1-point quadrature. To eliminate spurious zero-energy hourglass modes, the implementation in \texttt{simp\_engine\_3d.py} was updated to utilize full $2\times2\times2$ (8-point) Gauss quadrature, consistent with the construction of the elastic stiffness matrix.
    \item \textbf{Fix 2: ARPACK Out-Of-Memory (OOM) Crash in Buckling Solver.} The linear buckling solver originally invoked \texttt{scipy.sparse.linalg.eigsh} using the exact sparse matrix $\mathbf{K}_{\text{free}}$, prompting ARPACK to execute a memory-intensive SuperLU factorization (\texttt{splu}). To eliminate this OOM bottleneck, the solver was refactored to use shift-invert mode (\texttt{sigma=1e-6}) via a custom \texttt{LinearOperator} that iteratively solves the shifted system $(\mathbf{K}_{\text{free}} + \sigma \mathbf{G}_{\text{free}})\mathbf{x} = \mathbf{b}$ using Preconditioned Conjugate Gradient (PCG) with a Jacobi preconditioner. The problem is solved as $\mathbf{K} \mathbf{x} = \mu (-\mathbf{G}) \mathbf{x}$ in buckling mode to correctly find the smallest positive buckling load multiplier.
    \item \textbf{Fix 3: Remote Code Execution (RCE) Vulnerability in Dashboard.} The \texttt{eval()} calls used to parse mathematical expressions for nodal forces and boundary boxes in \texttt{app.py} presented a severe RCE vulnerability. These were successfully replaced with \texttt{sympy.sympify} to safely and robustly parse geometric user inputs.
    \item \textbf{Fix 4: Aggressive Heaviside Projection Ramp.} The exponent-based continuation ramp for the Heaviside projection parameter $\beta$ originally escalated too aggressively ($2^{\lfloor \text{it}/4 \rfloor}$), causing the optimizer to freeze in suboptimal chunky local minima. This was replaced with a gentler, linearly increasing ramp ($\beta = \min(32.0, 1.0 + \text{it} / 10.0)$), affording the algorithm sufficient iteration cycles to discover highly slender truss-like ties.
\end{itemize}

\section{Sprint 2: Multi-Physics Extension for Thermo-Elastic Optimization}

In order to evolve the static compliance topology optimization codebase into a comprehensive multi-physics engine, Sprint 2 introduces the theoretical and computational infrastructure required to concurrently solve the linear elastic equilibrium and steady-state heat conduction problems.

\subsection{Modular Assembly Architecture}
To accommodate multiple Partial Differential Equations (PDEs), the monolithic \texttt{solve()} method was refactored into an object-oriented modular framework. Element stiffness and load assemblies are dynamically encapsulated into targeted subroutines (\texttt{assemble\_elastic\_}\allowbreak\texttt{stiffness}, \texttt{assemble\_geometric\_}\allowbreak\texttt{stiffness}, \texttt{assemble\_thermal\_}\allowbreak\texttt{conductivity}, and \texttt{assemble\_thermal\_}\allowbreak\texttt{load\_vector}). This isolates the physics definitions from the iterative projection, filtering, and optimizer loop. 

\subsection{3D H8 Heat Conduction and Thermal Expansion Elements}
The steady-state heat conduction equation over the spatial domain is discretized using the same 3D H8 (eight-node hexahedral) isoparametric element. Each node possesses a single thermal degree of freedom (temperature $T$), leading to an $8 \times 8$ elemental thermal conductivity matrix $\mathbf{k}_{\text{th}, 0}$. For an isotropic thermal conductivity $k_{\text{th}}$, the constitutive matrix is strictly diagonal:
\begin{equation}
\mathbf{D}_{\text{th}} = k_{\text{th}} \mathbf{I}_3 = \begin{bmatrix}
k_{\text{th}} & 0 & 0 \\
0 & k_{\text{th}} & 0 \\
0 & 0 & k_{\text{th}}
\end{bmatrix}
\end{equation}
The thermal gradient matrix $\mathbf{B}_{\text{th}} \in \mathbb{R}^{3 \times 8}$ is constructed from the natural derivatives of the standard trilinear shape functions, and $\mathbf{k}_{\text{th}, 0}$ is synthesized via full $2 \times 2 \times 2$ Gauss-Legendre quadrature:
\begin{equation}
\mathbf{k}_{\text{th}, 0} = \int_{-1}^1 \int_{-1}^1 \int_{-1}^1 \mathbf{B}_{\text{th}}^T \mathbf{D}_{\text{th}} \mathbf{B}_{\text{th}} \det(\mathbf{J}) \, d\xi \, d\eta \, d\zeta
\end{equation}

Subsequently, local temperature changes $\Delta T_e$ generate an initial thermal strain vector $\boldsymbol{\varepsilon}_{\text{th}} = [\alpha_{\text{th}}\Delta T_e, \alpha_{\text{th}}\Delta T_e, \alpha_{\text{th}}\Delta T_e, 0, 0, 0]^T$, where $\alpha_{\text{th}}$ is the coefficient of thermal expansion. The equivalent thermal load vector $\mathbf{f}_{\text{th},0} \in \mathbb{R}^{24}$ mapping the thermal expansion into nodal forces is evaluated as:
\begin{equation}
\mathbf{f}_{\text{th},0} = \int_{\Omega_e} \mathbf{B}^T \mathbf{D} \boldsymbol{\varepsilon}_{\text{th}} \, d\Omega
\end{equation}

\subsection{RAMP Interpolation for Thermo-Elastic Coupling}
\label{sec:ramp}
Standard SIMP penalization ($E(x) = x^p E_0$, with $p=3$) proves mathematically ill-suited for thermo-elastic topology optimization \cite{pedersen2010, rodrigues1995}. Because SIMP artificially drives the elastic stiffness toward zero as $\hat{x}_e^3$, whereas a linear thermal expansion model scales with $\hat{x}_e$, the local strain ratio scales as $\varepsilon_{\text{th}} \sim F_{\text{th}} / K \sim \hat{x}_e / \hat{x}_e^3 = \hat{x}_e^{-2} \to \infty$ in low-density void regions ($\hat{x}_e \to 0$). Consequently, intermediate and void elements exhibit pathological, unbounded fictitious displacements that distort the global compliance calculation and destabilize the numerical optimizer.

Furthermore, if thermal expansion loads were penalized with standard SIMP ($p=3$), its derivative with respect to density would vanish identically at zero density ($\frac{dE}{d\hat{x}_e}(0) = 0$). This derivative nullification creates artificial local minima and causes the optimizer to freeze in void regions where sensitivity signals are completely extinguished.

To resolve both pathologies simultaneously and preserve well-posedness, the \textit{Rational Approximation of Material Properties} (RAMP) interpolation scheme \cite{bendsoe2003, pedersen2010} is introduced explicitly for the coupled thermal expansion load vector:
\begin{equation}
E_{\text{th}}(\hat{x}_e) = E_{\min} + \frac{\hat{x}_e}{1 + q_{\text{RAMP}} (1 - \hat{x}_e)} (E_0 - E_{\min})
\end{equation}
where $q_{\text{RAMP}} = 8.0$ is the concavity curvature parameter, $E_0$ is the solid Young's modulus, and $E_{\min} = 10^{-6} E_0$ is the minimum stiffness floor preventing matrix singularity.

Differentiating $E_{\text{th}}(\hat{x}_e)$ with respect to pseudo-density yields:
\begin{equation}
\frac{dE_{\text{th}}}{d\hat{x}_e} = \frac{1 + q_{\text{RAMP}}}{\left(1 + q_{\text{RAMP}}(1 - \hat{x}_e)\right)^2} (E_0 - E_{\min})
\end{equation}
The RAMP formulation exhibits two vital mathematical properties:
\begin{enumerate}
    \item \textbf{Solid state ($\hat{x}_e = 1$):} $E_{\text{th}}(1) = E_0$, with steep sensitivity slope:
    \begin{equation}
    \frac{dE_{\text{th}}}{d\hat{x}_e}(1) = (1 + q_{\text{RAMP}})(E_0 - E_{\min}) = 9(E_0 - E_{\min})
    \end{equation}
    which strongly penalizes intermediate material states and promotes crisp $0-1$ discrete topologies.
    \item \textbf{Void state ($\hat{x}_e = 0$):} $E_{\text{th}}(0) = E_{\min}$, while retaining a strictly positive, non-zero derivative:
    \begin{equation}
    \frac{dE_{\text{th}}}{d\hat{x}_e}(0) = \frac{E_0 - E_{\min}}{1 + q_{\text{RAMP}}} = \frac{1}{9}(E_0 - E_{\min}) > 0
    \end{equation}
\end{enumerate}
Because the derivative remains non-zero at $\hat{x}_e = 0$, void elements receive finite sensitivity gradients throughout optimization, preventing numerical stagnation while guaranteeing physical compatibility between mechanical stiffness and thermal expansion.

\subsection{Coupled Thermo-Elastic Formulation and Adjoint Sensitivity Analysis}
\label{sec:thermoelastic}
The steady-state multi-physics problem couples heat conduction with linear elasticity under thermal expansion.

\subsubsection{Governing Coupled State Equations}
The forward physical system is governed by two sequential boundary-value problems:
\begin{enumerate}
    \item \textbf{Steady-state heat conduction:} Over the 3D computational domain $\Omega$, the nodal temperature field $\mathbf{T} \in \mathbb{R}^{N_{\text{nodes}}}$ satisfies:
    \begin{equation}
    \mathbf{K}_{\text{th}}(\hat{\mathbf{x}}) \mathbf{T} = \mathbf{Q}
    \end{equation}
    subject to prescribed Dirichlet temperature boundary conditions $T = T_{\text{fixed}}$ on boundary $\Gamma_T \subset \partial\Omega$. The global thermal conductivity matrix $\mathbf{K}_{\text{th}}(\hat{\mathbf{x}})$ is assembled from the elemental matrices $\mathbf{k}_{\text{th}, 0}$ via SIMP penalization ($p_{\text{th}} = 3.0$):
    \begin{equation}
    \mathbf{K}_{\text{th}}(\hat{\mathbf{x}}) = \sum_{e=1}^m \left[ k_{\min} + \hat{x}_e^{p_{\text{th}}} (1 - k_{\min}) \right] \mathbf{k}_{\text{th}, 0}^e
    \end{equation}
    with $k_{\min} = 10^{-6}$.
    \item \textbf{Elemental temperature rise:} For each hexahedral element $e$, the operating temperature change $\Delta T_e$ relative to the reference stress-free temperature $T_{\text{ref}}$ is calculated from the nodal temperature vector $\mathbf{T}_e \in \mathbb{R}^8$:
    \begin{equation}
    \Delta T_e = \bar{T}_e - T_{\text{ref}} = \frac{1}{8} \sum_{i=1}^8 T_{e, i} - T_{\text{ref}}
    \end{equation}
    \item \textbf{Global thermal load vector assembly:} The elemental thermal expansion generates an equivalent global nodal load vector $\mathbf{F}_{\text{th}}(\hat{\mathbf{x}}, \mathbf{T}) \in \mathbb{R}^{3 N_{\text{nodes}}}$ interpolated via RAMP:
    \begin{equation}
    \mathbf{F}_{\text{th}}(\hat{\mathbf{x}}, \mathbf{T}) = \sum_{e=1}^m E_{\text{th}}(\hat{x}_e) \Delta T_e \mathbf{f}_{\text{th}, 0}^e
    \end{equation}
    where $\mathbf{f}_{\text{th}, 0}^e \in \mathbb{R}^{24}$ is the pre-integrated unit thermal expansion vector for the H8 element, satisfying self-equilibrium ($\sum_{i=1}^8 \mathbf{f}_i = \mathbf{0}$).
    \item \textbf{Linear elastic static equilibrium:} The coupled displacement field $\mathbf{U} \in \mathbb{R}^{3 N_{\text{nodes}}}$ satisfies:
    \begin{equation}
    \mathbf{K}(\hat{\mathbf{x}}) \mathbf{U} = \mathbf{F}_{\text{total}} = \mathbf{F}_{\text{ext}} + \mathbf{F}_{\text{th}}(\hat{\mathbf{x}}, \mathbf{T})
    \end{equation}
    subject to kinematic displacement boundary conditions $\mathbf{U} = \mathbf{0}$ on $\Gamma_U \subset \partial\Omega$.
\end{enumerate}

\subsubsection{Coupled Compliance Objective Function}
Structural compliance $C(\hat{\mathbf{x}})$, representing the total external work done by mechanical and thermal loads, is defined as:
\begin{equation}
C(\hat{\mathbf{x}}) = \mathbf{F}_{\text{total}}^T \mathbf{U} = \left( \mathbf{F}_{\text{ext}} + \mathbf{F}_{\text{th}}(\hat{\mathbf{x}}, \mathbf{T}) \right)^T \mathbf{U}
\end{equation}
Exploiting static equilibrium $\mathbf{K}(\hat{\mathbf{x}}) \mathbf{U} = \mathbf{F}_{\text{total}}$, this expression simplifies to the symmetric quadratic form:
\begin{equation}
C(\hat{\mathbf{x}}) = \mathbf{U}^T \mathbf{K}(\hat{\mathbf{x}}) \mathbf{U} = \sum_{e=1}^m E_K(\hat{x}_e) \mathbf{u}_e^T \mathbf{k}_0 \mathbf{u}_e
\end{equation}
This formulation preserves complete consistency with purely mechanical topology optimization: when thermal loads are absent ($\mathbf{F}_{\text{th}} = \mathbf{0}$), $C$ reduces identically to the classical elastic strain energy.

\subsubsection{Thermal Adjoint System Derivation}
Because the state vector $\mathbf{U}$ depends on the temperature field $\mathbf{T}$, and both depend implicitly on the design pseudo-densities $\hat{\mathbf{x}}$, direct differentiation of $C$ with respect to each $\hat{x}_e$ would require $m$ expensive sensitivity solves. To circumvent this, we formulate the augmented Lagrangian $\mathcal{L}$ incorporating both mechanical and thermal state equations:
\begin{equation}
\mathcal{L}(\mathbf{U}, \mathbf{T}, \boldsymbol{\lambda}_U, \boldsymbol{\lambda}_T, \hat{\mathbf{x}}) = \mathbf{U}^T \mathbf{K} \mathbf{U} - \boldsymbol{\lambda}_U^T \left( \mathbf{K} \mathbf{U} - \mathbf{F}_{\text{ext}} - \mathbf{F}_{\text{th}}(\hat{\mathbf{x}}, \mathbf{T}) \right) - \boldsymbol{\lambda}_T^T \left( \mathbf{K}_{\text{th}} \mathbf{T} - \mathbf{Q} \right)
\end{equation}
where $\boldsymbol{\lambda}_U \in \mathbb{R}^{3 N_{\text{nodes}}}$ and $\boldsymbol{\lambda}_T \in \mathbb{R}^{N_{\text{nodes}}}$ are the adjoint Lagrange multiplier vectors.

Enforcing stationarity of $\mathcal{L}$ with respect to the state variables:
\begin{enumerate}
    \item \textbf{Stationarity with respect to displacement $\mathbf{U}$:}
    \begin{equation}
    \frac{\partial \mathcal{L}}{\partial \mathbf{U}} = 2 \mathbf{K} \mathbf{U} - \mathbf{K} \boldsymbol{\lambda}_U = \mathbf{0} \implies \boldsymbol{\lambda}_U = 2 \mathbf{U}
    \end{equation}
    Hence, the mechanical adjoint problem is self-adjoint, requiring no separate mechanical adjoint solve.
    \item \textbf{Stationarity with respect to temperature $\mathbf{T}$:}
    \begin{equation}
    \frac{\partial \mathcal{L}}{\partial \mathbf{T}} = \left( \frac{\partial \mathbf{F}_{\text{th}}}{\partial \mathbf{T}} \right)^T \boldsymbol{\lambda}_U - \mathbf{K}_{\text{th}} \boldsymbol{\lambda}_T = 2 \left( \frac{\partial \mathbf{F}_{\text{th}}}{\partial \mathbf{T}} \right)^T \mathbf{U} - \mathbf{K}_{\text{th}} \boldsymbol{\lambda}_T = \mathbf{0}
    \end{equation}
    Defining the scaled thermal adjoint vector $\mathbf{P} = \frac{1}{2} \boldsymbol{\lambda}_T \in \mathbb{R}^{N_{\text{nodes}}}$, we obtain the linear thermal adjoint system:
    \begin{equation}
    \mathbf{K}_{\text{th}}(\hat{\mathbf{x}}) \mathbf{P} = \mathbf{F}_{\text{adj, th}} = \left( \frac{\partial \mathbf{F}_{\text{th}}}{\partial \mathbf{T}} \right)^T \mathbf{U}
    \end{equation}
    where the global thermal adjoint load vector $\mathbf{F}_{\text{adj, th}}$ is accumulated across all elements incident to global node $i$:
    \begin{equation}
    (\mathbf{F}_{\text{adj, th}})_i = \sum_{e \in \text{elem}(i)} \frac{1}{8} E_{\text{th}}(\hat{x}_e) \left( \mathbf{u}_e^T \mathbf{f}_{\text{th}, 0}^e \right)
    \end{equation}
    Solving this single linear system with thermal Dirichlet boundary conditions ($\mathbf{P} = \mathbf{0}$ on $\Gamma_T$) yields the adjoint temperature field $\mathbf{P}$.
\end{enumerate}

\subsubsection{Complete 3-Term Adjoint Sensitivity Analysis}
Taking the total derivative of the Lagrangian $\mathcal{L}$ with respect to the element pseudo-density $\hat{x}_e$:
\begin{equation}
\frac{\partial C}{\partial \hat{x}_e} = \frac{\partial \mathcal{L}}{\partial \hat{x}_e} = \mathbf{u}_e^T \frac{\partial \mathbf{k}_e}{\partial \hat{x}_e} \mathbf{u}_e - \boldsymbol{\lambda}_{U, e}^T \left( \frac{\partial \mathbf{k}_e}{\partial \hat{x}_e} \mathbf{u}_e - \frac{\partial \mathbf{f}_{\text{th}, e}}{\partial \hat{x}_e} \right) - \boldsymbol{\lambda}_{T, e}^T \left( \frac{\partial \mathbf{k}_{\text{th}, e}}{\partial \hat{x}_e} \mathbf{T}_e \right)
\end{equation}
Substituting the adjoint solutions $\boldsymbol{\lambda}_{U, e} = 2 \mathbf{u}_e$ and $\boldsymbol{\lambda}_{T, e} = 2 \mathbf{p}_e$:
\begin{equation}
\frac{\partial C}{\partial \hat{x}_e} = \mathbf{u}_e^T \frac{\partial \mathbf{k}_e}{\partial \hat{x}_e} \mathbf{u}_e - 2 \mathbf{u}_e^T \frac{\partial \mathbf{k}_e}{\partial \hat{x}_e} \mathbf{u}_e + 2 \mathbf{u}_e^T \frac{\partial \mathbf{f}_{\text{th}, e}}{\partial \hat{x}_e} - 2 \mathbf{p}_e^T \frac{\partial \mathbf{k}_{\text{th}, e}}{\partial \hat{x}_e} \mathbf{T}_e
\end{equation}
Combining like terms yields the exact, closed-form three-term adjoint sensitivity formula:
\begin{equation}
\frac{\partial C}{\partial \hat{x}_e} = \underbrace{- \mathbf{u}_e^T \frac{\partial \mathbf{k}_e}{\partial \hat{x}_e} \mathbf{u}_e}_{\text{Term 1: Elastic Strain Energy}} + \underbrace{2 \mathbf{u}_e^T \frac{\partial \mathbf{f}_{\text{th}, e}}{\partial \hat{x}_e}}_{\text{Term 2: RAMP Thermal Expansion}} - \underbrace{2 \mathbf{p}_e^T \frac{\partial \mathbf{k}_{\text{th}, e}}{\partial \hat{x}_e} \mathbf{T}_e}_{\text{Term 3: Thermal Conductivity Adjoint}}
\end{equation}
where each term is evaluated analytically as:
\begin{align}
\text{Term 1} &= - p_K \hat{x}_e^{p_K - 1} (E_0 - E_{\min}) \left( \mathbf{u}_e^T \mathbf{k}_0 \mathbf{u}_e \right) \\[4pt]
\text{Term 2} &= 2 \left[ \frac{1 + q_{\text{RAMP}}}{\left(1 + q_{\text{RAMP}}(1 - \hat{x}_e)\right)^2} (E_0 - E_{\min}) \right] \Delta T_e \left( \mathbf{u}_e^T \mathbf{f}_{\text{th}, 0}^e \right) \\[4pt]
\text{Term 3} &= - 2 p_{\text{th}} \hat{x}_e^{p_{\text{th}} - 1} (1 - k_{\min}) \left( \mathbf{p}_e^T \mathbf{k}_{\text{th}, 0} \mathbf{T}_e \right)
\end{align}

\subsubsection{Physical Interpretation, Sign Reversal, and Numerical Validation}
In purely mechanical optimization, Term 1 is strictly negative ($\partial C / \partial \hat{x}_e \le 0$), guaranteeing that adding material always stiffens the structure and decreases compliance. 

In thermo-elastic systems, however, Term 2 introduces a fundamentally different physical mechanism: when a structure expands thermally against rigid kinematic constraints, adding material increases the magnitude of thermal expansion forces, which may deform the structure further and increase compliance ($\partial C / \partial \hat{x}_e > 0$). This sign reversal violates the basic monotonic assumptions of the classical Optimality Criteria (OC) update, causing it to fail or oscillate uncontrollably.

The analytical three-term adjoint formulation was rigorously validated against central finite differences with perturbation step size $\epsilon = 10^{-6}$:
\begin{itemize}
    \item \textbf{Coupled Heat Conduction + Elasticity:} Relative sensitivity error between analytical adjoint and finite differences was verified at:
    \begin{equation}
    \frac{|\partial C_{\text{adjoint}} - \partial C_{\text{FD}}|}{|\partial C_{\text{FD}}|} = 2.14 \times 10^{-6}
    \end{equation}
    \item \textbf{Uniform Temperature Field ($\Delta T = 25^\circ\text{C}$):} With constant temperature, $\frac{\partial \mathbf{T}}{\partial \hat{x}_e} = \mathbf{0} \implies \mathbf{P} = \mathbf{0}$, reducing the sensitivity to Term 1 + Term 2. Relative error was verified at $1.55 \times 10^{-6}$.
\end{itemize}

\subsection{Method of Moving Asymptotes (MMA) Optimizer Integration}
\label{sec:mma}
To accommodate non-monotonic sensitivities, sign reversals in thermo-elastic compliance, and general multi-constrained formulations, the heuristic Optimality Criteria (OC) scheme is superseded by Svanberg's Method of Moving Asymptotes (MMA) \cite{svanberg1987}.

\subsubsection{Svanberg's Rational Approximating Subproblem}
Consider the general non-linear programming problem:
\begin{equation}
\begin{aligned}
\min_{\mathbf{x} \in \mathbb{R}^n} \quad & f_0(\mathbf{x}) \\
\text{subject to} \quad & f_i(\mathbf{x}) \le 0, \quad i = 1, \dots, m \\
& \alpha_j \le x_j \le \beta_j, \quad j = 1, \dots, n
\end{aligned}
\end{equation}
At iteration $k$, MMA approximates the objective and constraint functions by strictly convex, separable rational approximations $\tilde{f}_i^{(k)}(\mathbf{x})$:
\begin{equation}
\tilde{f}_i^{(k)}(\mathbf{x}) = r_i^{(k)} + \sum_{j=1}^n \left( \frac{p_{ij}^{(k)}}{U_j^{(k)} - x_j} + \frac{q_{ij}^{(k)}}{x_j - L_j^{(k)}} \right), \quad i = 0, 1, \dots, m
\end{equation}
where $L_j^{(k)}$ and $U_j^{(k)}$ are lower and upper moving asymptotes satisfying $L_j^{(k)} < \alpha_j^{(k)} \le x_j^{(k)} \le \beta_j^{(k)} < U_j^{(k)}$.
The rational coefficients $p_{ij}^{(k)} \ge 0$ and $q_{ij}^{(k)} \ge 0$ are determined by matching first derivatives:
\begin{equation}
p_{ij}^{(k)} = \begin{cases}
\left( U_j^{(k)} - x_j^{(k)} \right)^2 \frac{\partial f_i}{\partial x_j}\left(\mathbf{x}^{(k)}\right), & \text{if } \frac{\partial f_i}{\partial x_j}\left(\mathbf{x}^{(k)}\right) > 0 \\[6pt]
0, & \text{if } \frac{\partial f_i}{\partial x_j}\left(\mathbf{x}^{(k)}\right) \le 0
\end{cases}
\end{equation}
\begin{equation}
q_{ij}^{(k)} = \begin{cases}
0, & \text{if } \frac{\partial f_i}{\partial x_j}\left(\mathbf{x}^{(k)}\right) \ge 0 \\[6pt]
- \left( x_j^{(k)} - L_j^{(k)} \right)^2 \frac{\partial f_i}{\partial x_j}\left(\mathbf{x}^{(k)}\right), & \text{if } \frac{\partial f_i}{\partial x_j}\left(\mathbf{x}^{(k)}\right) < 0
\end{cases}
\end{equation}
and the scalar offset $r_i^{(k)}$ is chosen such that $\tilde{f}_i^{(k)}(\mathbf{x}^{(k)}) = f_i(\mathbf{x}^{(k)})$:
\begin{equation}
r_i^{(k)} = f_i\left(\mathbf{x}^{(k)}\right) - \sum_{j=1}^n \left( \frac{p_{ij}^{(k)}}{U_j^{(k)} - x_j^{(k)}} + \frac{q_{ij}^{(k)}}{x_j^{(k)} - L_j^{(k)}} \right)
\end{equation}
Because the second derivatives are strictly positive:
\begin{equation}
\frac{\partial^2 \tilde{f}_i^{(k)}}{\partial x_j^2} = \frac{2 p_{ij}^{(k)}}{(U_j^{(k)} - x_j)^3} + \frac{2 q_{ij}^{(k)}}{(x_j - L_j^{(k)})^3} > 0
\end{equation}
each subproblem is strictly convex and separable, guaranteeing a unique global minimum.

\subsubsection{Moving Asymptotes Update Heuristics}
The positions of the asymptotes control the conservatism and step size of each subproblem. When asymptotes are tight, curvature is high, restraining the design update; when asymptotes are distant, the model approaches a linear approximation. Svanberg's adaptive heuristic updates the asymptotes based on the history of design moves:

For the initial two iterations ($k=0, 1$):
\begin{equation}
L_j^{(k)} = x_j^{(k)} - s_0 (x_{j, \max} - x_{j, \min}), \quad U_j^{(k)} = x_j^{(k)} + s_0 (x_{j, \max} - x_{j, \min})
\end{equation}
with default scale $s_0 = 0.5$.

For subsequent iterations ($k \ge 2$), the movement history of each variable is tracked:
\begin{align}
L_j^{(k)} &= x_j^{(k)} - \gamma_j^{(k)} \left( x_j^{(k-1)} - L_j^{(k-1)} \right) \\
U_j^{(k)} &= x_j^{(k)} + \gamma_j^{(k)} \left( U_j^{(k-1)} - x_j^{(k-1)} \right)
\end{align}
where the scaling multiplier $\gamma_j^{(k)}$ adapts to oscillation or monotonicity:
\begin{equation}
\gamma_j^{(k)} = \begin{cases}
0.7, & \text{if } \left( x_j^{(k)} - x_j^{(k-1)} \right)\left( x_j^{(k-1)} - x_j^{(k-2)} \right) < 0 \quad (\text{oscillation: contract}) \\[4pt]
1.2, & \text{if } \left( x_j^{(k)} - x_j^{(k-1)} \right)\left( x_j^{(k-1)} - x_j^{(k-2)} \right) > 0 \quad (\text{monotonic: expand}) \\[4pt]
1.0, & \text{if } \left( x_j^{(k)} - x_j^{(k-1)} \right)\left( x_j^{(k-1)} - x_j^{(k-2)} \right) = 0
\end{cases}
\end{equation}
Safety bounds prevent the asymptotes from collapsing too close to the current design or escaping to infinity:
\begin{align}
x_j^{(k)} - 100(x_{j, \max} - x_{j, \min}) &\le L_j^{(k)} \le x_j^{(k)} - 0.01(x_{j, \max} - x_{j, \min}) \\
x_j^{(k)} + 0.01(x_{j, \max} - x_{j, \min}) &\le U_j^{(k)} \le x_j^{(k)} + 100(x_{j, \max} - x_{j, \min})
\end{align}

\subsubsection{Dual Subproblem Solving}
Because the MMA subproblem is strictly convex and separable, it is solved efficiently in the dual space with respect to the Lagrange multipliers $\boldsymbol{\lambda} = [\lambda_1, \dots, \lambda_m]^T \in \mathbb{R}^m_{\ge 0}$. The Lagrangian of the approximating subproblem is:
\begin{equation}
\mathcal{L}_{\text{MMA}}(\mathbf{x}, \boldsymbol{\lambda}) = \tilde{f}_0^{(k)}(\mathbf{x}) + \sum_{i=1}^m \lambda_i \tilde{f}_i^{(k)}(\mathbf{x})
\end{equation}
Minimizing $\mathcal{L}_{\text{MMA}}$ over the box constraints $\alpha_j \le x_j \le \beta_j$ produces the explicit, analytical primal-dual mapping $x_j^*(\boldsymbol{\lambda})$:
\begin{equation}
x_j^*(\boldsymbol{\lambda}) = \min\left( \beta_j, \, \max\left( \alpha_j, \, \frac{\sqrt{P_j(\boldsymbol{\lambda})} L_j^{(k)} + \sqrt{Q_j(\boldsymbol{\lambda})} U_j^{(k)}}{\sqrt{P_j(\boldsymbol{\lambda})} + \sqrt{Q_j(\boldsymbol{\lambda})}} \right) \right)
\end{equation}
where:
\begin{equation}
P_j(\boldsymbol{\lambda}) = p_{0j}^{(k)} + \sum_{i=1}^m \lambda_i p_{ij}^{(k)}, \quad Q_j(\boldsymbol{\lambda}) = q_{0j}^{(k)} + \sum_{i=1}^m \lambda_i q_{ij}^{(k)}
\end{equation}
The dual objective function $W(\boldsymbol{\lambda}) = \mathcal{L}_{\text{MMA}}(\mathbf{x}^*(\boldsymbol{\lambda}), \boldsymbol{\lambda})$ is concave and continuously differentiable, with gradient $\nabla W(\boldsymbol{\lambda}) = [\tilde{f}_1^{(k)}(\mathbf{x}^*(\boldsymbol{\lambda})), \dots, \tilde{f}_m^{(k)}(\mathbf{x}^*(\boldsymbol{\lambda}))]^T$. Because the number of constraints $m$ is typically very small ($m \ll n$, e.g., $m=1$ for volume conservation while $n \approx 10^4 - 10^6$), solving $\max_{\boldsymbol{\lambda} \ge \mathbf{0}} W(\boldsymbol{\lambda})$ via Newton's method or SQP requires negligible computation.

\subsubsection{NLopt Integration Architecture and Safeguards}
In the 3D topology optimization solver, MMA is implemented via the C++ backend of NLopt through its Python bindings (\texttt{nlopt.LD\_MMA}) \cite{johnson2014}. The software architecture incorporates four critical safeguards:
\begin{enumerate}
    \item \textbf{State-Caching Evaluator (\texttt{\_TOStateEvaluator}):} NLopt queries objective and constraint functions via separate callback routines. Evaluating the objective followed by the volume constraint within the same iteration would trigger duplicate $O(N)$ FEM solves. The caching evaluator caches $(\mathbf{x}, \mathbf{U}, \mathbf{T}, C, \nabla C, \mathbf{V}, \nabla \mathbf{V})$. When consecutive callbacks evaluate an identical design vector ($\|\mathbf{x} - \mathbf{x}_{\text{cached}}\|_\infty < 10^{-14}$), the cached results are returned immediately, eliminating $50\%$ of FEA computational overhead.
    \item \textbf{In-Place Gradient Buffer Assignment:} The NLopt Python interface passes a pre-allocated NumPy array view \texttt{grad} to the callback. Rebinding the local variable (\texttt{grad = ...}) fails to update the C-level memory buffer; gradients must be assigned in-place (\texttt{grad[:] = ...}) to prevent silent zero-gradient optimization stalls.
    \item \textbf{Objective Normalization:} Raw 3D compliance values frequently span $10^5 - 10^7\,\text{J}$. Such magnitudes distort the internal asymptote scaling within NLopt. The objective function is normalized by the initial compliance $C_0 = C(\mathbf{x}^{(0)})$:
    \begin{equation}
    f_0(\mathbf{x}) = \frac{C(\mathbf{x})}{C_0}, \quad \nabla f_0(\mathbf{x}) = \frac{\nabla C(\mathbf{x})}{C_0}
    \end{equation}
    keeping objective values and sensitivities strictly well-scaled ($O(1)$).
    \item \textbf{Passive Domains via Bound Constraints:} Prescribed solid non-design regions ($\hat{x}_e \equiv 1$) and void keep-out zones ($\hat{x}_e \equiv 0$) are enforced directly through variable bounds via \texttt{set\_lower\_bounds} and \texttt{set\_upper\_bounds}, guaranteeing exact geometric compliance without penalty distortions.
\end{enumerate}

\subsection{Algebraic Multigrid (AMG) Preconditioner Integration}
\label{sec:amg}
To solve large 3D linear systems for elasticity ($\mathbf{K} \mathbf{U} = \mathbf{F}$) and heat conduction ($\mathbf{K}_{\text{th}} \mathbf{T} = \mathbf{Q}$) on commodity hardware (8GB RAM), the solver integrates an Algebraic Multigrid (AMG) preconditioner via \texttt{pyamg} \cite{bell2012}.

\subsubsection{Theoretical Motivation and Scaling Limits}
Direct sparse factorization methods (e.g. SuperLU, PARDISO) exhibit $O(N^{1.5}-N^2)$ memory scaling and $O(N^2-N^3)$ computational complexity in 3D, quickly exhausting 8GB of RAM for meshes exceeding $20 \times 10 \times 10$ elements. Conversely, iterative solvers with simple diagonal (Jacobi) preconditioning suffer from spectral degradation: the condition number of 3D elliptic operators scales as $\kappa(\mathbf{K}) = O(h^{-2}) = O(N^{2/3})$, demanding thousands of Conjugate Gradient iterations.

Smoothed Aggregation Algebraic Multigrid (SA-AMG) \cite{vanek1996} overcomes these limitations by constructing a hierarchy of coarse spaces purely from matrix entries, achieving optimal $O(N)$ computational complexity and mesh-independent convergence without requiring geometric grid coordinates.

\subsubsection{Smoothed Aggregation Theory (SA-AMG)}
Given the sparse Symmetric Positive Definite (SPD) system matrix $\mathbf{A} \in \mathbb{R}^{N \times N}$:
\begin{enumerate}
    \item \textbf{Strength of Connection:} Degrees of freedom $i$ and $j$ are strongly coupled if their off-diagonal coupling exceeds a relative threshold:
    \begin{equation}
    |A_{ij}| \ge \theta_{\text{AMG}} \sqrt{A_{ii} A_{jj}}
    \end{equation}
    where $\theta_{\text{AMG}} \approx 0.08$ is the strength parameter.
    \item \textbf{Aggregation:} The nodes of the adjacency graph are partitioned into disjoint neighborhood aggregates $\{\mathcal{A}_k\}_{k=1}^{N_c}$.
    \item \textbf{Near-Nullspace Modes and Tentative Prolongator:} Standard point smoothers (e.g. symmetric Gauss-Seidel or damped Jacobi) efficiently eliminate high-energy oscillatory errors, leaving smooth, low-energy error modes in the near-nullspace of $\mathbf{A}$. These modes must be accurately represented on coarse levels via candidate vectors $\mathbf{B} \in \mathbb{R}^{N \times d_{\text{ns}}}$:
    \begin{itemize}
        \item \textbf{Steady-state heat conduction ($\mathbf{K}_{\text{th}}$):} $d_{\text{ns}} = 1$. The near-nullspace mode is the constant temperature shift vector:
        \begin{equation}
        \mathbf{B}_{\text{th}} = \mathbf{1} \in \mathbb{R}^N
        \end{equation}
        consistent with the single zero eigenvalue of $\mathbf{k}_{\text{th}, 0}$.
        \item \textbf{3D Linear Elasticity ($\mathbf{K}$):} $d_{\text{ns}} = 6$. The near-nullspace modes comprise 3 rigid-body translations and 3 infinitesimal rigid-body rotations:
        \begin{equation}
        \mathbf{B}_{\text{elas}} = \begin{bmatrix}
        \mathbf{1} & \mathbf{0} & \mathbf{0} & \mathbf{0} & \mathbf{z} & -\mathbf{y} \\
        \mathbf{0} & \mathbf{1} & \mathbf{0} & -\mathbf{z} & \mathbf{0} & \mathbf{x} \\
        \mathbf{0} & \mathbf{0} & \mathbf{1} & \mathbf{y} & -\mathbf{x} & \mathbf{0}
        \end{bmatrix}
        \end{equation}
    \end{itemize}
    On each aggregate $\mathcal{A}_k$, local QR factorization of $\mathbf{B}|_{\mathcal{A}_k}$ defines the tentative prolongator $\mathbf{P}_{\text{tent}} \in \mathbb{R}^{N \times (N_c \cdot d_{\text{ns}})}$, guaranteeing exact partition of unity and coarse representation of near-nullspace modes ($\mathbf{P}_{\text{tent}} \mathbf{B}_c = \mathbf{B}$).
    \item \textbf{Prolongator Smoothing and Galerkin Operator:} To remove high-frequency energy from the coarse basis, $\mathbf{P}_{\text{tent}}$ is smoothed with a damped Jacobi filter:
    \begin{equation}
    \mathbf{P} = \left( \mathbf{I} - \omega \mathbf{D}^{-1} \mathbf{A} \right) \mathbf{P}_{\text{tent}}
    \end{equation}
    where $\mathbf{D} = \text{diag}(\mathbf{A})$ and $\omega = \frac{4}{3 \lambda_{\max}(\mathbf{D}^{-1}\mathbf{A})}$. The coarse-grid Galerkin operator is constructed via projection:
    \begin{equation}
    \mathbf{A}_c = \mathbf{P}^T \mathbf{A} \mathbf{P}
    \end{equation}
    with restriction operator $\mathbf{R} = \mathbf{P}^T$.
\end{enumerate}

\subsubsection{Preconditioned Conjugate Gradient (PCG) Acceleration}
Rather than deploying AMG as a standalone stationary solver, a single AMG $V$-cycle ($\text{cycle}='V'$) serves as the preconditioner $\mathbf{M}^{-1}$ within the Krylov Conjugate Gradient algorithm (\texttt{scipy.}\allowbreak\texttt{sparse.}\allowbreak\texttt{linalg.}\allowbreak\texttt{cg}). 

Because the multilevel cycle dampens error components across all spatial scales simultaneously, the spectrum of $\mathbf{M}^{-1}\mathbf{A}$ is tightly clustered, bounding the condition number independently of mesh size:
\begin{equation}
\kappa\left( \mathbf{M}^{-1} \mathbf{A} \right) \le C_{\text{AMG}} = O(1)
\end{equation}
As a result, PCG converges in a mesh-independent number of iterations (typically 15--35 iterations), delivering a 5--10$\times$ speedup over Jacobi-PCG while consuming negligible memory.

\subsubsection{Defensive Safeguards for Topology Optimization}
Topology optimization produces extreme material stiffness contrast ($E_0 / E_{\min} = 10^6$), creating ill-conditioned linear systems that can destabilize standard multigrid setups. The engine implements four defensive numerical safeguards:
\begin{enumerate}
    \item \textbf{Dirichlet Submatrix Partitioning:} Constrained boundary degrees of freedom are explicitly partitioned prior to constructing the multigrid hierarchy:
    \begin{equation}
    \mathbf{A}_{\text{free}} = \mathbf{A}[\text{free\_dofs}, \text{free\_dofs}], \quad \mathbf{b}_{\text{free}} = \mathbf{b}[\text{free\_dofs}]
    \end{equation}
    Purging constrained boundary rows prevents uncoupled boundary nodes from corrupting aggregate formation and ensures that $\mathbf{A}_{\text{free}}$ is strictly Symmetric Positive Definite.
    \item \textbf{SVD Pseudoinverse Coarse Grid Solver (\texttt{coarse\_solver='pinv'}):} On the coarsest multigrid level, default direct sparse LU factorizations (\texttt{splu}) frequently crash or produce float overflow when porous coarse topologies generate near-singular matrices. Using an SVD-based Moore-Penrose pseudoinverse (\texttt{pinv}) truncates near-zero singular values, guaranteeing unconditional robustness against singularity.
    \item \textbf{Strict \texttt{float64} Precision Enforcement:} All sparse matrix arrays (\texttt{data}, \texttt{indices}, \texttt{indptr}) and load vectors are strictly enforced in \texttt{np.float64} precision. Reduced precision (\texttt{float32}) causes rapid loss of Krylov orthogonality in high-contrast void-solid meshes.
    \item \textbf{Three-Tier Fallback Hierarchy:} The \texttt{solve\_linear\_system} method implements an automated recovery cascade:
    \begin{equation}
    \text{SA-AMG PCG} \quad \xrightarrow{\text{on failure}} \quad \text{Jacobi-PCG} \quad \xrightarrow{\text{on failure}} \quad \text{SuperLU Direct Solve}
    \end{equation}
    ensuring that intermediate numerical anomalies never raise unhandled exceptions or crash the optimization run.
\end{enumerate}

\begin{thebibliography}{99}
\bibitem{ferrari2020}
Ferrari, F., \& Sigmund, O. (2020). Revisiting topology optimization with buckling constraints. \textit{Structural and Multidisciplinary Optimization}, 61(4), 1401--1415.

\bibitem{sigmund2001}
Sigmund, O. (2001). A 99 line topology optimization code written in Matlab. \textit{Structural and Multidisciplinary Optimization}, 21(2), 120--127.

\bibitem{marler2004}
Marler, R. T., \& Arora, J. S. (2004). Survey of multi-objective optimization methods for engineering. \textit{Structural and Multidisciplinary Optimization}, 26(6), 369--395.

\bibitem{sigmund1997}
Sigmund, O. (1997). On the design of compliant mechanisms using topology optimization. \textit{Mechanics of Structures and Machines}, 25(4), 493--524.

\bibitem{neves1995}
Neves, M. M., Rodrigues, H., \& Guedes, J. M. (1995). Generalized topology design of structures with a buckling load criterion. \textit{Structural Optimization}, 10(2), 71--78.

\bibitem{bendsoe2003}
Bendsøe, M. P., \& Sigmund, O. (2003). \textit{Topology Optimization: Theory, Methods, and Applications}. Springer Science \& Business Media.

\bibitem{ferrari2021}
Ferrari, F., Sigmund, O., \& Guest, J. K. (2021). Topology optimization with linearized buckling criteria in 250 lines of Matlab. \textit{Structural and Multidisciplinary Optimization}, 63(6), 3045--3066.

\bibitem{wang2011}
Wang, F., Lazarov, B. S., \& Sigmund, O. (2011). On projection methods, convergence and robust formulations in topology optimization. \textit{Structural and Multidisciplinary Optimization}, 43(6), 767--784.

\bibitem{das1997}
Das, I., \& Dennis, J. E. (1997). A closer look at drawbacks of minimizing weighted sums of objectives for Pareto set generation in multicriteria optimization problems. \textit{Structural Optimization}, 14(1), 63--69.

\bibitem{geoffrion1968}
Geoffrion, A. M. (1968). Proper efficiency and the theory of vector maximization. \textit{Journal of Mathematical Analysis and Applications}, 22(3), 618--630.

\bibitem{taubin1995}
Taubin, G. (1995). Curve and surface smoothing without shrinkage. In \textit{Proceedings of IEEE International Conference on Computer Vision} (pp. 852--857). IEEE.

\bibitem{svanberg1987}
Svanberg, K. (1987). The method of moving asymptotes---a new method for structural optimization. \textit{International Journal for Numerical Methods in Engineering}, 24(2), 359--373.

\bibitem{vanek1996}
Van\v{e}k, P., Mandel, J., \& Brezina, M. (1996). Algebraic multigrid by smoothed aggregation for second and fourth order elliptic problems. \textit{Computing}, 56(3), 179--196.

\bibitem{bell2012}
Bell, W. N., Olson, L. N., \& Schroder, J. B. (2012). PyAMG: Algebraic multigrid solvers in Python. \textit{ACM Transactions on Mathematical Software (TOMS)}.

\bibitem{pedersen2010}
Pedersen, N. L., \& Pedersen, P. (2010). Design of thermoelastic structures using topology optimization. \textit{Structural and Multidisciplinary Optimization}, 42(5), 701--713.

\bibitem{rodrigues1995}
Rodrigues, H., \& Fernandes, P. (1995). Topology optimal design of thermoelastic structures using a homogenization method. \textit{Computer Methods in Applied Mechanics and Engineering}, 125(1--4), 1--20.

\bibitem{johnson2014}
Johnson, S. G. (2014). The NLopt nonlinear-optimization package. \url{https://github.com/stevengj/nlopt}.
\end{thebibliography}

\end{document}
